\scanpage{125}
\mspage{352}

is in the principal class, i.e. we will have
\deleted{\unsure{(19)}}
\[
\alpha\rho(\alpha)^{-1}=\alpha_0\rho(\alpha_0)^{-1}\qquad\text{for }\alpha_0\in\Pi\text{ suitable}
\]
In fact, we can write $\alpha=\alpha_0\rho^i$\inserted{, with $\alpha_0\in\Pi$, $i\in\mathbb Z$}, whence $\rho(\alpha)^{-1}=\rho^{-i}\rho(\alpha_0)^{-1}$ whence (19). More generally, we will have, \deleted{\unsure{if $\alpha,\alpha_0$}}
\[
\alpha=\alpha_0\rho^i\qquad(\deleted{\unsure{\alpha\in\widetilde\Pi}}\alpha_0\in\Pi,\ i\in\mathbb Z)
\]
the relation
\deleted{\unsure{(19)}}
\[
\begin{aligned}
\alpha g\rho(\alpha^{-1})
&=\alpha_0\underset{\deleted{\unsure{\rho^i(g)}}}{\underbrace{\rho^i g\rho^{-i}}}\unsure{\rho(\alpha_0^{-1})|\rho(\alpha_0)^{-1}}\\
&=\deleted{\unsure{\alpha_0\rho^i g\rho(\alpha_0\rho^i)^{-1}}}\inserted{\alpha_0\rho^i(g)\rho(\alpha_0)^{-1}}
\end{aligned}
\]
\unsure{with|for}
\[
\rho^i(g)=\beta g\rho(\beta)^{-1}\qquad
\beta=(g\rho(g)\cdots\rho^{i-1}(g))^{-1}
\]
thus
\[
\left\{\begin{aligned}
\alpha g\rho(\alpha^{-1})&=(\alpha_0\beta)g\rho(\alpha_0\beta)^{-1}\\
\beta&=(g\rho(g)\cdots\rho^{i-1}(g))^{-1}
\end{aligned}\right. \tag{19}
\]
(for $\alpha=\alpha_0\rho^i$, $i\in\unsure{\mathbb N|\mathbb N^*}$).

Of course, the representation (18) of a solution of equation (8), \deleted{\unsure{is not}} when we no longer assume $\alpha\in\Pi$ but only $\alpha\in\widetilde\Pi$, is no longer one-to-one, since for $\alpha=\rho$ for example we find \deleted{\unsure{the element 1}} $g=\rho\rho^{-1}=1$.

\deleted{Let us suppose \unsure{in} addition}\inserted{Let us suppose} more generally that we have a group
\deleted{\unsure{(20)}}
\[
\unsure{\mathcal E|\varepsilon}\supset\widetilde\Pi
\]
(as in the \deleted{\unsure{group}} case where $\Pi=\Pi_{0,3}$, $\widetilde\Pi$ inverse image of $\{1,\dot\rho,\dot\rho^2\}\subset\mathfrak S_3$ in $\mathfrak S_{0,3}$ extension of $\mathcal T_{0,3}$ by $\Pi_{0,3}$), in which $\Pi$ is normal, so that \deleted{$\unsure{\mathcal E|\varepsilon}$ operates by conjugation}\inserted{we have an extension}

\scanpage{126}

\deleted{on the \unsure{set} of $\rho'$ in $\Pi$ \unsure{such|such} that}

\[
(20)\qquad 1\longrightarrow\Pi\longrightarrow\unsure{\mathcal E|\varepsilon}\longrightarrow G\longrightarrow1
\]

\margin{N.B. the extension (3) is the inverse image of (20) under the inclusion $\mathbb Z/n\mathbb Z\hookrightarrow G$}

\unsure{with}
\[
(21)\qquad \mathbb Z/n\mathbb Z\hookrightarrow G,\qquad
1\bmod n\longmapsto\dot\rho\in G
\]
\inserted{(element of order $n$)}
$\widetilde\Pi$ being identified with the inverse image of $\mathbb Z/n\mathbb Z$ in $\unsure{\mathcal E|\varepsilon}$.

We cannot say in general, if $\rho'$ is an \unsure{element} of $\unsure{\mathcal E|\varepsilon}$ above $\dot\rho$ (thus $\rho'\in\widetilde{\unsure{\mathcal E|\varepsilon}}$).
(even if it satisfies $\rho'{}^n=1$), that for $\forall$ $\alpha\in\unsure{\mathcal E|\varepsilon}$, $\operatorname{int}(\alpha)(\rho')$ is still above $\dot\rho$:
for this to be so, it \unsure{is necessary and sufficient} that $\dot\alpha\in G$ centralize $\dot\rho$. So let $\unsure{\mathcal E|\varepsilon}_\rho$
\[
(22)\qquad
\unsure{\mathcal E|\varepsilon}_\rho=\{\alpha\in\unsure{\mathcal E|\varepsilon}\mid[\dot\alpha,\dot\rho]=1\}
\]
it is a subgroup of $\unsure{\mathcal E|\varepsilon}$ containing $\widetilde\Pi$, which
\deleted{\unsure{for $\forall$}}\inserted{\unsure{itself}} operates on the set of \deleted{\unsure{$\rho'$}} $\rho'$ above $\dot\rho$,
\inserted{in particular on \deleted{\illegible{1}} the set of \illegible{1} \unsure{which}}
\unsure{satisfying} (8), \deleted{and it} by the same operation
\[
(23)\qquad
\alpha\rho'_g\alpha^{-1}=\rho'_{\alpha g\rho(\alpha)^{-1}}
\]
so that we can also consider this \unsure{operation} as given by
\[
(24)\qquad
\left\{
\begin{array}{l}
\Theta_{\rho,\unsure{\mathcal E|\varepsilon}_\rho}(\alpha)=\Theta(\alpha):
g\longmapsto\alpha g\rho(\alpha)^{-1}\\
(\alpha\in\unsure{\mathcal E|\varepsilon}_\rho).
\end{array}
\right.
\]

This time, an operation $\Theta(\alpha)$ may permute the $\Pi$-conjugacy classes of the \deleted{\unsure{various}} solutions, of (8) in $\Pi$.
\unsure{We|It} \deleted{\illegible{1}} will have on $\unsure{\mathcal E|\varepsilon}_\rho$

\scanpage{127}
\mspage{353}

an extension structure
\[
(25)\qquad
\begin{array}{ccccccccc}
1&\longrightarrow&\Pi&\longrightarrow&\unsure{\mathcal E|\varepsilon}_\rho&\longrightarrow&G^\rho&\longrightarrow&1\\
&&&&&&\Vert\\
&&&&&&\multicolumn{3}{l}{\operatorname{Centr}_G(\rho)=\{g\in G\mid\rho g\unsure{g|\rho}^{-1}=g\}}\\
&&&&&&\multicolumn{3}{l}{\phantom{\operatorname{Centr}_G(\rho)}=\{g\in G\mid[\rho,g]=1\}}
\end{array}
\]
and $\unsure{\mathcal E|\varepsilon}_\rho$ operates on the set of $\Pi$-conjugacy classes of splittings of (3), (or the $\Pi$-$\rho$-conjugacy classes of solutions of (8)) via an operation of $G^\rho$ on it.

Let $\deleted{\unsure{g}}\alpha\in\unsure{\mathcal E|\varepsilon}$, let us consider the element
\[
\deleted{\unsure{g}}\alpha\rho(\deleted{\unsure{g}}\alpha)^{-1}
=\alpha(\rho\alpha\rho^{-1})^{-1}
=\alpha\rho\alpha^{-1}\rho^{-1}=[\alpha,\rho]
\]
It is in $\Pi$ iff we have \deleted{\unsure{that}}
\[
[\dot\alpha,\dot\rho]=1\qquad\text{i.e.}\quad
\dot\alpha\in G^\rho\qquad\text{i.e.}\quad\alpha\in\unsure{\mathcal E|\varepsilon}_\rho.
\]

For a solution $g$ of equation (8), it may happen that we cannot put \unsure{it|it} in the \inserted{standard} form $\alpha\rho(\alpha)^{-1}$ with $\alpha\in\Pi$, but that we succeed with $\alpha\in\unsure{\mathcal E|\varepsilon}_\rho$.
For such an $\alpha$ to exist, it is necessary and sufficient that the $\Pi\rho$-conjugacy class \inserted{$[g]_\rho$} of $g$ (among the solutions of (8)) be in the orbit of the trivial conjugacy class \inserted{$[1]_\rho$}, for the action under consideration of $G^\rho$ on this set of conjugacy classes; \deleted{and}\inserted{then} the $\alpha\in\unsure{\mathcal E|\varepsilon}_\rho$ which \unsure{satisfy}
\[
(26)\qquad g=\alpha\rho(\alpha)^{-1}
\]

\scanpage{128}

\deleted{are such that}
\deleted{satisfy}
\deleted{\unsure{are some} \illegible{1}}
\inserted{\unsure{satisfy} the relation}
\[
(27)\qquad\dot\alpha[1]_\rho=[g]_\rho,
\]
and for $\forall\ \dot\alpha\in G^\rho$ satisfying this relation (27), there exists at least one representative $\alpha\in\unsure{\mathcal E|\varepsilon}_\rho$ satisfying (6) (which is \unsure{pure} tautology).

Under what condition do we have \uline{uniqueness} \unsure{of the|of the} $\alpha\in\unsure{\mathcal E|\varepsilon}_\rho$ satisfying (26)?
The set of these $\alpha$ is a right coset under the \unsure{subgroup} of $\unsure{\mathcal E|\varepsilon}_\rho$ \unsure{formed} of the $\alpha$ such that \deleted{\illegible{1}}
\[
\alpha\rho(\alpha)^{-1}=1\qquad\text{i.e.}\quad
\alpha=\rho(\alpha)\qquad\text{i.e.}\quad[\alpha,\rho]=1,
\]
this group is none other than \deleted{$\unsure{\mathcal E|\varepsilon}$}
\[
(28)\qquad\unsure{\mathcal E|\varepsilon}^\rho=\operatorname{Centr}_{\unsure{\mathcal E|\varepsilon}}(\rho)\subset\unsure{\mathcal E|\varepsilon}_\rho.
\]
\unsure{which} group contains $\rho$ \unsure{hence} the \inserted{\unsure{sub-}}group $\simeq\mathbb Z/n\mathbb Z$ \inserted{of $\unsure{\mathcal E|\varepsilon}$} generated by $\rho$ -- so it is never trivial -- i.e. we never have uniqueness (except in the trivial case where \deleted{\unsure{$\rho$}} $n=1$ i.e. $\rho=1$, \deleted{\illegible{1}} but then $\unsure{\mathcal E|\varepsilon}^\rho=\unsure{\mathcal E|\varepsilon}$ and we would need $\unsure{\mathcal E|\varepsilon}=1$ \dots !)

But if we assume $\Pi^\rho=\{1\}$ i.e.
\[
(29)\qquad\unsure{\mathcal E|\varepsilon}^\rho\cap\Pi=\{1\}
\]
\deleted{\unsure{then}}
\[
\unsure{\mathcal E|\varepsilon}^\rho_{\deleted{\unsure{\rho}}}
\underset{\uline{\text{injective}}}{\overset{\mathrm{can}}{\hookrightarrow}}
G^\rho\subset G,
\]
then the $\alpha\in\unsure{\mathcal E|\varepsilon}^\rho$ solutions of (26) are in one-to-one correspon-

\scanpage{129}
\mspage{354}

dence with their images \inserted{$\dot\alpha$} in $G$, i.e. with the transporter \inserted{in $G$} from $[1]_\rho$ \unsure{into|to} $[g]_\rho$.
If we choose an element of this transporter, we have chosen by \unsure{that very fact} an $\alpha$, solution of (26).

\[
\boxed{\begin{gathered}
\text{In the case of }\quad
\Pi=\Pi_{\unsure{0,3}}\inserted{=\Pi_1(U_{0,3})},\qquad
\unsure{\mathcal E|\varepsilon}=\mathfrak S_{0,3}
\inserted{=\Pi_1(U_{0,3},\mathcal T_{0,3})},
\quad\deleted{\unsure{G=\mathcal T_0}}\\
G=\mathcal T_{0,3}\simeq\mathfrak S_3\times\mathbb Z/2\mathbb Z
\end{gathered}}
\]
with $\rho,\dot\rho$ as usual, the set of conjugacy classes of the \deleted{\unsure{solutions}}\inserted{splittings} of (3) i.e. of the solutions of (8), \unsure{is identified} with $U_{0,3}^{\dot\rho}$, which has two elements, \unsure{the} two ``poles'' $-\overline j,-j$ -- one of which, chosen as base point, corresponds to the trivial class.

\deleted{\unsure{Here the} \illegible{1} $G^\rho\simeq\unsure{\mathbb Z/2\mathbb Z}\{1,\tau\}\simeq\mathbb Z$ is the group}
\inserted{Here}
\[
G^\rho=\unsure{\mathcal T|\mathcal G}_{0,3}^{\rho}
\simeq\mathbb Z/3\mathbb Z\times\mathbb Z/2\mathbb Z
\]
is the group generated by $\rho$ and $\tau$, and $\rho$ operates \unsure{trivially|trivially} on the set of the two classes, while $\tau$ permutes them (i.e. $\mathbb Z/\unsure{2|n}\mathbb Z$ \inserted{generated by $\tau$} operates on it in a simply transitive way).
Thus if we \unsure{set}
\[
\begin{aligned}
(30)\qquad\unsure{\mathcal E|\varepsilon}'
&=\inserted{\unsure{\widetilde\Pi|\Pi}_{0,3}=}\Pi\cdot\{1,\tau_\infty\}\\
&=\text{inverse image of }\{1,\tau_\infty\}\subset G\\
&\phantom{={}}\text{in }\unsure{\mathcal E|\varepsilon}=\mathfrak S_{0,3}
\end{aligned}
\]
\margin{(\deleted{\illegible{1}} unoriented \unsure{weighted} triangulated cartographic \unsure{group} \dots}
we see that
\[
(31)\qquad\unsure{\mathcal E|\varepsilon}'{}^\rho=\{1\},
\]
\deleted{(immediate consequence of the relation $\Pi^\rho=(1)$)}
and the set of solutions in \uline{$\Pi_{0,3}$} of the equation
\[
(32)\qquad g\rho(g)\rho^2(g)=1
\]

\scanpage{130}

is in 1-1 correspondence with the set of elements of $\widetilde\Pi_{0,3}=\unsure{\mathcal E|\varepsilon}'$ (the unoriented weighted \inserted{triangulated} cartographic group, having as generators the elements $\tau_0,\tau_1,\tau_\infty$ related by the relations $\tau_0^2=\tau_1^2=\tau_\infty^2=1$ and no more, with $l_0=\tau_\infty\tau_1$, $l_1=\tau_0\tau_\infty$, $l_\infty=\tau_1\tau_0$ \dots)
It should still be likewise when we place ourselves in the groups $\Pi_{0,3}^{\wedge}$ and $\widetilde\Pi_{0,3}^{\wedge}=\widetilde\Pi_{0,3}\cdot\{1,\tau_\infty\}$, if we admit (something to be checked later)
\[
(33)\qquad(\widetilde\Pi_{0,3}^{\wedge})^\rho=\{1\}
\]
\inserted{and that the elements of finite order of $\widehat{\mathfrak S}_{0,3}$ are conjugates of elements of finite order of $\mathfrak S_{0,3}$ \dots}

Let us note that here we have
\[
(34)\qquad\tau_\infty\rho(\tau_\infty)^{-1}
=\tau_\infty\rho(\tau_\infty)=l_1^{-1},
\qquad\deleted{\unsure{\text{in}}}
\]
in virtue of the formulas (23) of \S\ 42
\[
\rho_\infty\overset{\unsure{\mathrm{dfn}|\text{definition}}}{=}
\tau_\infty\rho\tau_\infty^{-1}
=l_1^{-1}\rho=\rho l_0^{-1}.
\]

The situation is even simpler for the \unsure{element|element} $\sigma_\infty$ of $\mathfrak S_{0,3}$, operating on $\Pi$ by an automorphism of order $n=2$, for here every solution of the \unsure{corresponding} equation
\[
(35)\qquad\deleted{\unsure{(36)}}\quad
\deleted{\unsure{g}}h\,\sigma_\infty(\deleted{\unsure{g}}h)=1
\qquad(h\in\Pi)
\]
can be put in the standard form
\[
(36)\qquad h=\beta\sigma_\infty(\beta)^{-1}
\qquad(\beta\in\Pi)
\]
and this in a \uline{unique way}. This \unsure{follows} indeed from the general theory \dots, taking into account that
\[
(37)\qquad U_{0,3}^{\dot\sigma_\infty}=\{\unsure{\tfrac12}\}
\]